\documentclass[11pt]{article}

\usepackage[utf8]{inputenc}
\usepackage[T1]{fontenc}
\usepackage{lmodern}
\usepackage{amsmath,amssymb,amsthm,mathtools}
\usepackage{stmaryrd}
\usepackage[margin=1in]{geometry}
\usepackage{enumitem}
\usepackage{xcolor}
\usepackage{booktabs}
\usepackage{hyperref}
\usepackage[nameinlink,capitalize]{cleveref}

\hypersetup{
  colorlinks=true,
  linkcolor=blue!60!black,
  citecolor=blue!60!black,
  urlcolor=blue!60!black
}

% theorem environments
\newtheorem{theorem}{Theorem}
\newtheorem{lemma}[theorem]{Lemma}
\newtheorem{proposition}[theorem]{Proposition}
\newtheorem{corollary}[theorem]{Corollary}
\newtheorem{question}[theorem]{Question}
\newtheorem{conjecture}[theorem]{Conjecture}
\theoremstyle{definition}
\newtheorem{definition}[theorem]{Definition}
\newtheorem{example}[theorem]{Example}
\newtheorem{remark}[theorem]{Remark}

% macros
\newcommand{\R}{\mathbb{R}}
\newcommand{\N}{\mathbb{N}}
\newcommand{\E}{\mathbb{E}}
\newcommand{\Prob}{\mathbb{P}}
\newcommand{\A}{\mathcal{A}}
\newcommand{\F}{\mathcal{F}}
\newcommand{\K}{\mathcal{K}}
\newcommand{\T}{\mathcal{T}}
\newcommand{\Homp}{\operatorname{Hom}^{+}}
\newcommand{\Homalg}{\operatorname{Hom}}
\newcommand{\Csem}{\mathcal{C}_{\mathrm{sem}}}
\newcommand{\Ext}{\operatorname{Ext}}
\newcommand{\supp}{\operatorname{supp}}
\newcommand{\im}{\operatorname{im}}
\newcommand{\brak}[1]{\left\langle #1\right\rangle}
\newcommand{\avg}[2]{\left\llbracket #1 \right\rrbracket_{#2}}
\newcommand{\q}{\mathfrak{q}}
\newcommand{\bvec}{\mathbf{m}}
\newcommand{\eps}{\varepsilon}

\title{Forbidden-Subgraph Reasoning via Random Extensions:\\
       A Blow-Up Equivalence in Flag Algebras}
\author{}
\date{\today}

\begin{document}
\maketitle

\begin{abstract}
In Razborov's flag-algebra calculus there are two natural ways to express
that an algebra element~$f$ is asymptotically non-negative under a forbidden
subgraph constraint. The first descends to the quotient algebra in which the
forbidden configurations have been axiomatised away; the second remains in
the ambient theory, restricts the unlabelled limit homomorphism to be
supported on the constrained class, and asks for almost-sure non-negativity
under Razborov's canonical random extension to a labelled type.
We prove that the two formulations are equivalent for every \emph{hereditary
class closed under independent blow-ups}.  This includes all clique-free,
multipartite, and many other natural forbidden families, in both undirected
graphs and any relational theory admitting an independent blow-up
construction (digraphs, $r$-uniform hypergraphs, $\ldots$).

The easy direction is a clean consequence of Razborov's support transfer
through downward averaging.  The hard direction is proved by a
\emph{planted blow-up} construction: a labelled flag in the constrained
theory is realised, with positive probability, inside the random extension
of a weighted independent blow-up whose unlabelled limit is supported on the
constrained class.  We compare this with the alternative proof in the
unlabelled case via the (induced) graph removal lemma, and explain
precisely why each method works in a regime where the other does not.
\end{abstract}

\tableofcontents

\section{Introduction}

\subsection{What this paper is about}

Razborov's flag algebras provide a uniform language for asymptotic
extremal combinatorics~\cite{Razborov2007}.  At the centre of the
formalism is a commutative associative algebra $\A^\sigma[\T]$ attached to
each finite first-order type~$\sigma$ in a universal relational theory~$\T$;
the elements of $\A^\sigma[\T]$ are formal limits of densities of
$\sigma$-labelled finite models.  A frequent and useful operation is to
\emph{forbid} certain induced configurations.  Concretely, if $\K$ is a
hereditary class of $\T$-models, the constrained theory $\T_\K$ gives rise
to a quotient algebra
\[
  \q_\sigma : \A^\sigma[\T] \twoheadrightarrow \A^\sigma[\T_\K]
\]
which kills every flag whose underlying model lies outside~$\K$
(Razborov~\cite[\S2.3.3]{Razborov2007}).  The natural \emph{quotient
semantics} declares an element $f\in\A^\sigma[\T]$ to be non-negative
under the constraint if
\begin{equation}\label{eq:quotient}
  \psi(\q_\sigma f)\ge 0
  \qquad\text{for every }\psi\in\Homp(\A^\sigma[\T_\K],\R).
\end{equation}

In many calculations it is cleaner to keep working with the ambient
algebra $\A^\sigma[\T]$, never passing to the quotient.  Then the
constraint enters only through the choice of the unlabelled root
homomorphism: one fixes
$\phi_0\in\Homp(\A^0[\T],\R)$ that is \emph{supported on $\K$}
(i.e.\ $\phi_0(M)=0$ for every $M\notin\K$), and one looks at the
\emph{random labelled extension} $\phi^\sigma\sim\Ext_\sigma(\phi_0)$
guaranteed by Razborov's existence and uniqueness
theorem~\cite[Theorem~3.5]{Razborov2007}.  This leads to the
\emph{ensemble semantics}:
\begin{equation}\label{eq:ensemble}
  \Prob\bigl[\phi^\sigma(f)\ge 0\bigr]=1
  \qquad\text{for every $\K$-supported $\phi_0$ with $\phi_0(\brak\sigma)>0$.}
\end{equation}

Are \eqref{eq:quotient} and \eqref{eq:ensemble} the same statement?
For the unlabelled type ($\sigma=0$) the answer is essentially trivial:
the random extension is degenerate, so~\eqref{eq:ensemble} just says
that the unlabelled root is non-negative on $f$, which is~\eqref{eq:quotient}
read through the pull-back bijection between $\K$-supported
unlabelled-root homomorphisms and homomorphisms on the unlabelled
quotient algebra.  For labelled types, however, the question is
genuinely subtle: ensemble semantics quantifies only over those
$\K$-supported homomorphisms on $\A^\sigma[\T]$ that arise as random
extensions of unlabelled roots, and it is a priori unclear whether
these exhaust all $\K$-supported homomorphisms.  This sub-collection
question is in the spirit of the structural problems discussed at the
end of Razborov's paper~\cite[\S6]{Razborov2007} on the relationship
between syntactic statements in the calculus and the semantic objects
(ensembles of random homomorphisms) that test them.

\subsection{Our main result}

We give a clean sufficient condition under which the two semantics
coincide for every type $\sigma$:

\begin{quote}\em
  If $\K$ is closed under independent blow-ups, then quotient semantics
  and ensemble semantics agree for every labelled type and every element
  of $\A^\sigma[\T]$.
\end{quote}

This covers all $K_r$-free graph classes, hence in particular triangle-free
graphs, the setting that motivated this work; more generally it covers
multipartite classes, oriented graphs of bounded girth, and many other
natural forbidden families across graphs, digraphs, and hypergraphs.

The hard direction, ensemble~$\Rightarrow$~quotient, is the main
technical content.  The idea is short and we record it here: given a
finite labelled flag $F=(G,\theta)$ in the constrained theory $\T_\K$
that witnesses the failure of~\eqref{eq:quotient}, we build an
\emph{unlabelled} blow-up $B_N$ of $G$ in which a small but positive
fraction~$\rho$ of the vertices forms ``labelled clusters'' designed to
host random copies of the type~$\sigma$.  Each random embedding
$\widehat\theta:\sigma\to B_N$ that hits the planted clusters produces a
flag with essentially the same flag densities as $(G,\theta)$
(planted blow-up lemma).  The clone-closure hypothesis ensures
$B_N\in\K$, so the unlabelled limit of these blow-ups is $\K$-supported
and the bad flag lives, with positive measure, in the canonical random
extension of that limit, contradicting~\eqref{eq:ensemble}.

\subsection{A small example to fix ideas}

Take $\T=\T_{\mathrm{Graph}}$ and the triangle-free class
$\K=\mathrm{Forb}(K_3)$.  Triangle-free flag-algebra inequalities
typically take the shape ``$f\ge 0$ for triangle-free graphs'', where
$f\in\A^\sigma[\T_{\mathrm{Graph}}]$ is a labelled algebra element such
as the rooted edge $e\in\A^1$, the rooted square-degree
$\avg{e^2}{1}\in\A^0$, or any polynomial expression in these and other
rooted patterns.

\Cref{thm:main} guarantees two equivalent ways to certify such an
inequality:
\begin{itemize}[leftmargin=2em]
\item[(Q)]
verify $\psi(\q_\sigma f)\ge 0$ for every
$\psi\in\Homp(\A^\sigma[\T_\K],\R)$, working inside the quotient
algebra in which triangle-containing flags have been killed; or
\item[(E)]
verify, for every root homomorphism $\phi_0$ that puts mass~$0$ on
triangles and for almost every random extension
$\phi^\sigma\sim\Ext_\sigma(\phi_0)$, that $\phi^\sigma(f)\ge 0$,
working inside the ambient algebra
$\A^\sigma[\T_{\mathrm{Graph}}]$.
\end{itemize}
Form (E) lets one keep the full apparatus of the ambient calculus
(Cauchy--Schwarz~\cite[Theorem~3.14]{Razborov2007}, the differential
operators of~\cite[\S4.3]{Razborov2007}, all interpretations between
graph-theoretic objects, $\ldots$), with the triangle-free constraint
entering only through the choice of root.  Form (Q) is sometimes more
natural and lets one read off the inequality directly from the
quotient algebra structure.  Having both available is useful in
practice.

\subsection{Comparison with the removal-lemma approach}

There is a second, more analytic route to the equivalence, available
for the unlabelled type $\sigma=0$ and arbitrary hereditary $\K$:
every $\K$-supported limit homomorphism is a limit of finite $\K$-models
because the (induced) graph removal lemma of Alon, Fischer,
Krivelevich and Szegedy~\cite{AFKS2000}, extended by Alon and
Shapira~\cite{AlonShapira2008} to arbitrary hereditary properties,
lets us repair an arbitrary $H$-density-zero sequence into an
$H$-induced-free sequence with the same limit.  The two approaches are
complementary:

\begin{center}
\begin{tabular}{lcc}
\toprule
   & Removal-lemma route & Blow-up route (this paper)\\
\midrule
Hypothesis on $\K$           & Any hereditary class & Closed under blow-ups\\
Type $\sigma$                & Only $\sigma=0$       & Any non-degenerate $\sigma$\\
Quantitative ingredient      & Induced removal lemma & Planted blow-up\\
\bottomrule
\end{tabular}
\end{center}

Neither method currently delivers the strongest possible joint statement
(arbitrary hereditary $\K$, arbitrary type $\sigma$); we explain what
goes wrong in each combination in \Cref{sec:discussion}.

\subsection{Organisation}

\Cref{sec:setup} fixes notation and recalls the flag-algebra facts we
use as black boxes.  \Cref{sec:blowup} introduces the independent
blow-up axiomatically and gives the planted blow-up lemma.
\Cref{sec:proof} proves the main theorem.  \Cref{sec:discussion}
discusses generalisations, the relation with the removal-lemma route,
and open questions.

\section{Setup}\label{sec:setup}

\subsection{Theories admitting an independent blow-up}

Throughout, $\T$ is a universal first-order theory in a relational
language $L$ with equality but without function symbols or constants,
exactly as in~\cite[\S2]{Razborov2007}.  We assume $\T$ has at least one
model of every size and that every set of vertices of a model induces a
model (so the algebra $\A^\sigma[\T]$ is defined for every
non-degenerate type $\sigma$).  We write $V(M)$ for the vertex set of a
model~$M$ and $M\vert_S$ for the induced sub-model on $S\subseteq V(M)$.

\begin{definition}[Independent blow-up]\label{def:blow-up}
Let $M$ be a finite model of $\T$ and $\bvec=(m_v)_{v\in V(M)}$ a tuple
of positive integers.  An \emph{independent blow-up} of $M$ with
multiplicities $\bvec$ is a model $M^{\bvec}$ of $\T$ together with a
partition
\[
  V(M^{\bvec})=\bigsqcup_{v\in V(M)} C_v,
  \qquad |C_v|=m_v,
\]
called the \emph{cluster decomposition}, satisfying the
\emph{independence property}: for every subset
$S\subseteq V(M^{\bvec})$ containing at most one vertex from each cluster,
\[
  M^{\bvec}\bigl\vert_S \;\cong\; M\bigl\vert_{\pi(S)}
\]
via the projection $\pi\colon\bigsqcup_v C_v\to V(M)$ sending each
$w\in C_v$ to $v$.  The theory~$\T$ \emph{admits independent blow-ups} if
for every finite model~$M$ and every positive multiplicity vector
$\bvec$, an independent blow-up exists.
\end{definition}

\begin{example}\label{ex:graphs-and-friends}\hfill
\begin{enumerate}[label=\textup{(\alph*)}, leftmargin=2.5em]
\item For undirected simple graphs, the independent blow-up is the
classical construction: replace $v$ by an independent set $C_v$ of
size $m_v$ and join $C_v$ to $C_w$ by a complete bipartite graph
whenever $vw\in E(M)$.

\item For digraphs, replace $v$ by an independent set $C_v$ and add
\emph{all} arcs from $C_u$ to $C_w$ whenever $u\to w$ is an arc of $M$
(no arcs between different vertices of the same cluster, no arcs from a
vertex to itself).

\item For $r$-uniform hypergraphs, a set
$\{w_1,\ldots,w_r\}\subseteq V(M^{\bvec})$ is a hyperedge of $M^{\bvec}$
exactly when each $w_i$ lies in a distinct cluster $C_{v_i}$ and
$\{v_1,\ldots,v_r\}$ is a hyperedge of $M$.

\item More generally, for any universal theory in a relational language
whose axioms involve only distinct arguments, the recipe ``preserve
inter-cluster relations as in $M$, kill all intra-cluster relations''
defines an independent blow-up.
\end{enumerate}
\end{example}

\begin{definition}[Clone-closed class]\label{def:clone-closed}
A hereditary class $\K$ of finite $\T$-models is \emph{clone-closed} if
every independent blow-up of every $M\in\K$ also lies in $\K$.
\end{definition}

\begin{example}\label{ex:Kr-free-clone-closed}
For every $r\ge 2$, the class of $K_r$-free graphs is clone-closed.
Indeed, any clique in an independent blow-up meets each cluster in at
most one vertex (clusters are independent sets), so its projection is a
clique of the same size in the base graph.  In particular,
triangle-free graphs are clone-closed.

More generally:
\begin{itemize}[leftmargin=2em]
\item For graphs, the class $\mathrm{Forb}^{\mathrm{ind}}(\mathcal H)$ is
clone-closed whenever every $H\in\mathcal H$ is \emph{prime}, in the
sense that no two distinct non-adjacent vertices of $H$ have the same
neighbourhood inside $H$.  (Equivalently, $H$ has no \emph{false twins}.)
Any induced copy of such a prime $H$ in a blow-up has at most one vertex
per cluster (two vertices sharing a cluster would be non-adjacent false
twins in $H$) and therefore projects to a copy of $H$ in the base.
\item Multipartite classes ($r$-partite graphs, complete multipartite
graphs of bounded part count, etc.) are clone-closed, because each
cluster of a blow-up inherits the part of its base vertex.
\item For oriented graphs, the class of \emph{oriented graphs of girth
at least $g$} is clone-closed for every $g\ge 3$.  See
\Cref{sec:discussion} for the verification.
\end{itemize}
\end{example}

\begin{remark}[Why the hypothesis is needed]\label{rmk:why-hypothesis}
The clone-closure hypothesis fails for some hereditary classes.
For instance, $\mathrm{Forb}^{\mathrm{ind}}(P_3)$ consists of disjoint
unions of cliques, but the independent blow-up of $K_2$ with both
multiplicities equal to~$2$ is the four-cycle~$C_4$, which is not a
disjoint union of cliques.  The blow-up route in this paper genuinely
breaks down for such classes; the removal-lemma route still applies, but
only for the unlabelled type.  See \Cref{sec:discussion} for a
comparison.
\end{remark}

\subsection{Flag-algebra background, recalled briefly}

We use freely the syntactic and semantic apparatus of~\cite{Razborov2007}:
$\sigma$-flags $F=(M,\theta)$, the algebras $\A^\sigma[\T]$ with their
flag-density product, the homomorphism cone
$\Homp(\A^\sigma[\T],\R)$, the downward averaging operator
$\avg{\cdot}{\sigma,\eta}:\A^\sigma\to\A^{\sigma\vert_\eta}$, and the
constant $q_\sigma(F)\in[0,1]$ that records the labelling redundancy of a
flag.  When $\eta=0$ we write $\avg{\cdot}{\sigma}$.  We write
$\brak\sigma:=\avg{1_\sigma}\sigma\in\A^0[\T]$ for the unlabelled
counterpart of the identity flag, so by~\cite[Theorem~2.5\,(b)]{Razborov2007},
\[
  \brak\sigma = q_\sigma(1_\sigma)\cdot\sigma,
\]
where on the right we view $\sigma$ as an unlabelled model of $\T$.

The two black boxes we need are:

\begin{theorem}[Razborov's representation and extension theorems]\label{thm:razborov}
Let $\sigma$ be a non-degenerate type of a universal relational
theory~$\T$.
\begin{enumerate}[label=\textup{(\roman*)}, leftmargin=2.5em]
\item~\cite[Theorem~3.3]{Razborov2007}
Every convergent sequence of $\sigma$-flags
$F_n\in\F^\sigma$ with $|V(F_n)|\to\infty$ defines an element of
$\Homp(\A^\sigma,\R)$, and conversely every element of
$\Homp(\A^\sigma,\R)$ arises as the limit of such a sequence.
\item~\cite[Theorems~3.5,~3.12]{Razborov2007}
For every $\phi_0\in\Homp(\A^0[\T],\R)$ with $\phi_0(\brak\sigma)>0$,
there exists a unique probability measure $\Ext_\sigma(\phi_0)$ on
$\Homp(\A^\sigma[\T],\R)$ satisfying, for every $g\in\A^\sigma[\T]$,
\begin{equation}\label{eq:extension-expectation}
  \E_{\phi^\sigma\sim\Ext_\sigma(\phi_0)}\bigl[\phi^\sigma(g)\bigr]
  \;=\;
  \frac{\phi_0(\avg g\sigma)}{\phi_0(\brak\sigma)}.
\end{equation}
Moreover, this measure is the weak limit of the discrete random-root
distributions $P^\sigma_{M_n}$ obtained by sampling a labelled embedding
of $\sigma$ into $M_n$, for any sequence of $\T$-models $M_n$ with
$p^{M_n}\to\phi_0$.
\end{enumerate}
\end{theorem}

\subsection{\texorpdfstring{$\K$}{K}-supported homomorphisms}

\begin{definition}[$\K$-support]\label{def:K-support}
Let $\tau$ be a non-degenerate type realised in $\K$.  A homomorphism
$\varphi\in\Homp(\A^\tau[\T],\R)$ is \emph{$\K$-supported} if
$\varphi(F)=0$ for every $\tau$-flag $F$ whose underlying model is
outside $\K$.  Equivalently, $\varphi$ vanishes on $\ker(\q_\tau)$ and
hence factors uniquely through $\q_\tau$ to give an element of
$\Homp(\A^\tau[\T_\K],\R)$.
\end{definition}

The following two lemmas are routine; the second uses only the
expectation formula~\eqref{eq:extension-expectation} and not the
blow-up hypothesis.

\begin{lemma}[Pull-back identifies quotient with supported]\label{lem:quotient-support}
Composition with $\q_\tau$ is a bijection
\[
  \Homp(\A^\tau[\T_\K],\R)\;\xrightarrow{\;\sim\;}\;
  \bigl\{\,\varphi\in\Homp(\A^\tau[\T],\R)\,:\,\varphi\text{ is }\K\text{-supported}\bigr\}.
\]
\end{lemma}

\begin{proof}
The quotient kernel $\ker(\q_\tau)$ is the ideal generated by the
$\tau$-flags whose underlying model is outside $\K$.  Therefore every
homomorphism on the quotient pulls back to a $\K$-supported homomorphism,
and every $\K$-supported homomorphism vanishes on this kernel and
descends.
\end{proof}

\begin{lemma}[Random extensions preserve $\K$-support]\label{lem:support-passes}
Let $\phi_0\in\Homp(\A^0[\T],\R)$ be $\K$-supported and assume
$\phi_0(\brak\sigma)>0$.  If $\phi^\sigma\sim\Ext_\sigma(\phi_0)$, then
$\phi^\sigma$ is $\K$-supported almost surely.
\end{lemma}

\begin{proof}
Let $F$ be a $\sigma$-flag whose underlying unlabelled model $G$ is not
in $\K$.  By the definition of downward averaging
(\cite[\S2.2]{Razborov2007}),
\[
  \avg F \sigma \;=\; q_\sigma(F)\cdot G \;\in\; \A^0[\T],
\]
with $q_\sigma(F)\in[0,1]$.  Since $\phi_0$ is $\K$-supported and
$G\notin\K$, we have $\phi_0(\avg F\sigma)=0$, so by
\eqref{eq:extension-expectation},
\(
  \E[\phi^\sigma(F)]=0.
\)
As $\phi^\sigma(F)\ge 0$ pointwise (homomorphisms in $\Homp$ take flags
to non-negative reals), $\phi^\sigma(F)=0$ almost surely.  There are
countably many such bad flags, so a countable intersection of
probability-one events gives almost sure $\K$-support.
\end{proof}

\subsection{The main result}

\begin{theorem}[Forbidden-subgraph equivalence under blow-up closure]\label{thm:main}
Let $\T$ be a universal first-order relational theory admitting
independent blow-ups (\Cref{def:blow-up}).  Let $\K$ be a clone-closed
hereditary class of $\T$-models (\Cref{def:clone-closed}), and let
$\sigma$ be a non-degenerate type in $\T_\K$.  For every
$f\in\A^\sigma[\T]$, the following are equivalent.
\begin{enumerate}[label=\textup{(\alph*)}, leftmargin=2.5em]
  \item\label{item:quotient}\textup{(Quotient semantics.)}
  For every $\psi\in\Homp(\A^\sigma[\T_\K],\R)$,
  \(
     \psi(\q_\sigma f)\ge 0.
  \)
  \item\label{item:ensemble}\textup{(Ensemble semantics.)}
  For every $\K$-supported $\phi_0\in\Homp(\A^0[\T],\R)$ with
  $\phi_0(\brak\sigma)>0$ and its canonical random extension
  $\phi^\sigma\sim\Ext_\sigma(\phi_0)$,
  \(
     \Prob[\phi^\sigma(f)\ge 0]=1.
  \)
\end{enumerate}
\end{theorem}

\begin{corollary}[Clone-closed examples]\label{cor:cliques}
\Cref{thm:main} applies, in particular, to the following hereditary
classes (whose clone-closure is verified in
\Cref{ex:Kr-free-clone-closed} and \Cref{sec:discussion}):
\begin{enumerate}[label=\textup{(\alph*)}, leftmargin=2.5em]
\item undirected graphs forbidding any fixed clique $K_r$ \textup{($r\ge 2$)};
\item undirected $r$-partite graphs, for any fixed $r$;
\item undirected graphs $\mathrm{Forb}^{\mathrm{ind}}(\mathcal H)$ for
any family $\mathcal H$ of prime graphs (no false twins);
\item $r$-uniform hypergraphs forbidding any fixed complete
$r$-uniform sub-hypergraph $K_s^{(r)}$;
\item oriented graphs of girth at least $g$, for any $g\ge 3$.
\end{enumerate}
In each case, the equivalence holds for every non-degenerate type
$\sigma$ in the constrained theory and every $f\in\A^\sigma[\T]$.
\end{corollary}

\section{The planted blow-up lemma}\label{sec:blowup}

The hard direction of \Cref{thm:main} is driven by a single
quantitative idea: a labelled flag in a constrained theory can be
\emph{planted} into an independent blow-up so that, with positive
probability, the random extension of the unlabelled blow-up reproduces
the original labelled flag's densities up to arbitrarily small error.
We make this precise here.

\subsection{Weighted blow-ups}

Fix a $\sigma$-flag $F=(G,\theta)$ with $|\sigma|=k$ and $|V(G)|=n$, where
$n>k$.  Choose a parameter $0<\rho<1$.  Define weights
\[
  w_v=
  \begin{cases}
    \dfrac\rho k & \text{if }v=\theta(i)\text{ for some }i\in[k]\text{ and }k\ge 1,\\[1mm]
    \dfrac{1-\rho}{n-k} & \text{if }v\in V(G)\setminus\im\theta,
  \end{cases}
\]
which satisfy $\sum_v w_v=1$.  For an integer $N$, the
\emph{$(\rho,N)$-blow-up of $(G,\theta)$} is any independent blow-up
$B_N=B_N(G,\theta,\rho)$ of $G$ with cluster sizes $|C_v|$ as close to
$w_v\cdot N$ as the integer constraint allows.  Set
\[
  \operatorname{err}_N
  \;:=\;
  \sum_{v\in V(G)}\left\lvert\frac{|C_v|}{N}-w_v\right\rvert.
\]
A standard greedy rounding makes $\operatorname{err}_N\le n/N$, so
$\operatorname{err}_N\to 0$ as $N\to\infty$.

When $k=0$ (no labelled vertices to plant), there is no use for $\rho$:
take uniform weights $w_v=1/n$ and the same rounding.  We keep $\rho$ in
the notation for uniformity; the bounds below remain trivially true.

The cluster $C_{\theta(i)}$ associated to a labelled vertex is called a
\emph{labelled cluster}; the others are \emph{unlabelled clusters}.

\subsection{Planting a labelled embedding}

Given an unlabelled blow-up $B_N$, a labelled embedding
$\widehat\theta\colon[k]\to V(B_N)$ defining a $\sigma$-flag
$(B_N,\widehat\theta)$ is called \emph{planted} if
$\widehat\theta(i)\in C_{\theta(i)}$ for every $i\in[k]$.
By the independence property of the blow-up, any planted choice
$\widehat\theta$ is automatically a model embedding of $\sigma$ into
$B_N$: each $\widehat\theta(i)$ projects to $\theta(i)$, and the induced
sub-model on $\{\widehat\theta(1),\ldots,\widehat\theta(k)\}$ is
isomorphic to $\sigma$.

\subsection{The two key estimates}

\begin{lemma}[Planted approximation]\label{lem:planted}
Fix an integer $m\ge k$.  There exists a constant $C_m>0$ (depending only
on $m$ and the language $L$) such that the following holds.  For every
$\sigma$-flag $(G,\theta)$ with $|V(G)|=n>k$, every $(\rho,N)$-blow-up
$B_N=B_N(G,\theta,\rho)$, every planted labelled embedding
$\widehat\theta\colon[k]\to V(B_N)$, and every $\sigma$-flag $F_0$ with
at most $m$ vertices,
\begin{equation}\label{eq:planted-flag-bound}
  \bigl\lvert p(F_0,(B_N,\widehat\theta))-p(F_0,(G,\theta))\bigr\rvert
  \;\le\;
  C_m\!\left(\rho+\frac1{n-k}+\operatorname{err}_N\right).
\end{equation}
Consequently, for every finite linear combination $f\in\A^\sigma[\T]$ of
flags of size at most $m$ there is a constant $C_f>0$ with
\begin{equation}\label{eq:planted-f-bound}
  \bigl\lvert p(f,(B_N,\widehat\theta))-p(f,(G,\theta))\bigr\rvert
  \;\le\;
  C_f\!\left(\rho+\frac1{n-k}+\operatorname{err}_N\right).
\end{equation}
\end{lemma}

\begin{proof}
By linearity, \eqref{eq:planted-f-bound} follows from
\eqref{eq:planted-flag-bound}; we prove the latter.

Let $\ell=|V(F_0)|$; since $F_0$ is a $\sigma$-flag we automatically
have $\ell\ge k$, and by hypothesis $\ell\le m$.  Computing
$p(F_0,(B_N,\widehat\theta))$ amounts to sampling $\ell-k$ additional
vertices uniformly without replacement from
$V(B_N)\setminus\im\widehat\theta$ and asking for an isomorphism of
labelled flags.  Call such a sample \emph{good} if
\begin{enumerate}[label=\textup{(\roman*)}, leftmargin=2.5em]
  \item none of the sampled vertices lies in a labelled cluster, and
  \item no two sampled vertices lie in the same unlabelled cluster.
\end{enumerate}

We first estimate the probability of failure of each condition.
The total mass of labelled clusters in
$V(B_N)\setminus\im\widehat\theta$ is $\rho\pm O(\operatorname{err}_N)$,
so by a union bound the probability of (i) failing is at most
$(\ell-k)\bigl(\rho+O(\operatorname{err}_N)\bigr)$.  For each pair of
sample positions and each unlabelled cluster $C_v$, the probability
that both positions fall in $C_v$ is at most
$\bigl((1-\rho)/(n-k)\bigr)^{2}+O(\operatorname{err}_N/(n-k))$; summing
over $\binom{\ell-k}{2}$ pairs and $n-k$ unlabelled clusters gives
\[
  \Prob[\text{(ii) fails}]\;\le\;\frac{\binom{\ell-k}{2}}{n-k}+O_m(\operatorname{err}_N).
\]

On the good event, the projection
$\pi\bigl(\{\widehat\theta(1),\ldots,\widehat\theta(k)\}\cup S\bigr)$
of the sampled set $S$ together with the labelled vertices is, by the
independence property of the blow-up, an induced sub-model of $G$
isomorphic via $\pi$ to the corresponding induced sub-model of $B_N$.
Moreover, conditional on the good event, the projection of $S$ is a
uniformly random $(\ell-k)$-subset of $V(G)\setminus\im\theta$ up to
total variation $O_m(\operatorname{err}_N)$ caused by unequal cluster
sizes.  Hence on the good event the two density-computing
distributions agree up to total variation $O_m(\operatorname{err}_N)$.

Combining the bad-event mass with the good-event total-variation gap
gives a total deviation at most
$C_m\bigl(\rho+1/(n-k)+\operatorname{err}_N\bigr)$ for a constant $C_m$
depending only on $m$ and the language.
\end{proof}

\begin{lemma}[Positive mass of the planted event]\label{lem:positive-mass}
Set
\[
  c_{k,\rho}\;:=\;
  \begin{cases}
    1, & k=0,\\[1mm]
    \left(\dfrac{\rho}{2k}\right)^{\!k}, & k\ge 1.
  \end{cases}
\]
For every $\sigma$-flag $(G,\theta)$ and every sufficiently large $N$,
the discrete random root distribution $P^\sigma_{B_N}$ on
$\Homp(\A^\sigma,\R)$ (defined in~\cite[Definition~9]{Razborov2007} by
sampling a uniform labelled embedding $\widehat\theta\colon\sigma\to B_N$)
assigns probability at least $c_{k,\rho}$ to the event that
$\widehat\theta$ is planted.
\end{lemma}

\begin{proof}
For $k=0$ there is nothing to prove.  Assume $k\ge 1$.  Every ordered
choice of $\widehat\theta(i)\in C_{\theta(i)}$ is a model embedding of
$\sigma$ by the independence property.  For $N$ large enough that
$|C_{\theta(i)}|\ge\rho N/(2k)$ for each~$i$, the number of planted
ordered embeddings is at least $(\rho/(2k))^k\,N^k$.  The total number
of ordered injections of $[k]$ into $V(B_N)$, which upper bounds the
total number of $\sigma$-embeddings, is at most $N^k$.  Dividing gives
the claim.
\end{proof}

\section{Proof of the main theorem}\label{sec:proof}

\subsection{Quotient implies ensemble}

Assume~\ref{item:quotient}.  Let $\phi_0\in\Homp(\A^0[\T],\R)$ be
$\K$-supported with $\phi_0(\brak\sigma)>0$, and sample
$\phi^\sigma\sim\Ext_\sigma(\phi_0)$.  By \Cref{lem:support-passes},
$\phi^\sigma$ is $\K$-supported almost surely.  By
\Cref{lem:quotient-support}, on this almost-sure event $\phi^\sigma$
descends to $\overline{\phi}{}^\sigma\in\Homp(\A^\sigma[\T_\K],\R)$
with $\overline{\phi}{}^\sigma\circ\q_\sigma=\phi^\sigma$.
Applying~\ref{item:quotient} to $\overline{\phi}{}^\sigma$,
\[
  \phi^\sigma(f)
  \;=\;\overline{\phi}{}^\sigma(\q_\sigma f)
  \;\ge 0
  \qquad\text{almost surely.}
\]
This proves~\ref{item:ensemble}.

\subsection{Ensemble implies quotient}

The case $k=|\sigma|=0$ is the trivial pull-back equivalence
discussed in \Cref{sec:discussion}, so we assume $k\ge 1$
throughout this subsection.

Assume~\ref{item:ensemble}; we prove~\ref{item:quotient} by
contradiction.  Suppose there exists
$\psi_0\in\Homp(\A^\sigma[\T_\K],\R)$ with $\psi_0(\q_\sigma f)<0$.
Pull $\psi_0$ back along $\q_\sigma$ and write
$\psi:=\psi_0\circ\q_\sigma\in\Homp(\A^\sigma[\T],\R)$; then $\psi$ is
$\K$-supported and $\psi(f)=\psi_0(\q_\sigma f)<0$.  Choose $\alpha>0$
with
\begin{equation}\label{eq:negative-margin}
  \psi(f)\le-4\alpha.
\end{equation}

\paragraph{Step 1: Approximation in the constrained theory.}
By \Cref{thm:razborov}(i) applied to the theory $\T_\K$, there is an
increasing convergent sequence of $\sigma$-flags
\[
  F_n=(G_n,\theta_n)\in\F^\sigma[\T_\K],\qquad|V(G_n)|\to\infty,
\]
with $p^{F_n}\to\psi$ on $\A^\sigma[\T_\K]$.  By the construction of
$\q_\sigma$, this same sequence satisfies $p^{F_n}\to\psi$ on
$\A^\sigma[\T]$ as well, since any $\sigma$-flag whose underlying model
is outside $\K$ has flag density zero in every $F_n$ and is sent to zero
by $\psi$.  After passing to a tail,
\begin{equation}\label{eq:finite-negative}
  p(f,F_n)\le-3\alpha
  \qquad\text{for every }n.
\end{equation}

\paragraph{Step 2: Building $\K$-supported blow-ups.}
Let $m$ be the maximum size of any flag occurring in a fixed linear
expansion of $f$, and let $C_f$ be the constant from
\Cref{lem:planted}.  Choose $0<\rho<1$ with $C_f\rho\le\alpha/2$.  For
each $n$, choose $N_n$ so large that the rounding error
$C_f\operatorname{err}_{N_n}<\alpha/4$ in the $(\rho,N_n)$-blow-up
$B_n:=B_{N_n}(G_n,\theta_n,\rho)$, and so that $|V(B_n)|\to\infty$.
Since $G_n\in\K$ and $\K$ is clone-closed, each $B_n$ is in $\K$.

By compactness, pass to a subsequence (not relabelled) along which the
unlabelled graphs $B_n$ converge to some
$\phi_0\in\Homp(\A^0[\T],\R)$.  Because every $B_n$ is in $\K$ and $\K$
is hereditary, $p(M,B_n)=0$ for every $M\notin\K$, so taking the limit
gives $\phi_0(M)=0$; that is, $\phi_0$ is $\K$-supported.

\paragraph{Step 3: The type density is non-trivial.}
By the proof of \Cref{lem:positive-mass}, the blow-up $B_n$ admits at
least $(\rho/(2k))^k\cdot|V(B_n)|^k$ ordered model embeddings of
$\sigma$.  Each induced unlabelled copy of $\sigma$ in $B_n$ supports
exactly $|\mathrm{Aut}(\sigma)|$ such embeddings, so the unlabelled flag
density satisfies
\[
  p(\sigma, B_n)
  \;=\; \frac{\#\{\text{ordered embeddings }\sigma\to B_n\}}
              {|\mathrm{Aut}(\sigma)|\cdot\binom{|V(B_n)|}{k}}
  \;\ge\; \frac{(\rho/(2k))^k\cdot|V(B_n)|^k}{|\mathrm{Aut}(\sigma)|\cdot\binom{|V(B_n)|}{k}}
  \;\ge\; \frac{k!}{|\mathrm{Aut}(\sigma)|}\left(\frac{\rho}{2k}\right)^{\!k},
\]
where the last inequality uses
$|V(B_n)|^k/\binom{|V(B_n)|}{k}\ge k!$.  This lower bound is positive and
independent of $n$.  Hence
\[
  \phi_0(\brak\sigma)
  \;=\; q_\sigma(1_\sigma)\cdot \phi_0(\sigma)
  \;=\; q_\sigma(1_\sigma)\cdot\lim_n p(\sigma,B_n)
  \;>\;0,
\]
so~\ref{item:ensemble} applies to the random extension
$\phi^\sigma\sim\Ext_\sigma(\phi_0)$.

\paragraph{Step 4: Locating the contradiction.}
Let $P_n:=P^\sigma_{B_n}$ be the discrete random-root distribution on
the compact product space $X_\sigma:=[0,1]^{\F^\sigma}$.  By
\Cref{lem:positive-mass}, for $n$ large,
\[
  P_n\bigl(\bigl\{\widehat\theta\text{ planted}\bigr\}\bigr)
  \;\ge\; c_{k,\rho}.
\]
On the planted event, \Cref{lem:planted} together with
\eqref{eq:finite-negative} and the choices
$C_f\rho\le\alpha/2$,
$C_f\operatorname{err}_{N_n}<\alpha/4$ yield, for all $n$ with
$|V(G_n)|-k>4C_f/\alpha$,
\[
  p(f,(B_n,\widehat\theta))
  \;\le\; p(f,F_n)+C_f\!\left(\rho+\tfrac1{|V(G_n)|-k}+\operatorname{err}_{N_n}\right)
  \;\le\; -3\alpha+\tfrac\alpha2+\tfrac\alpha4+\tfrac\alpha4
  \;=\; -2\alpha.
\]
Let
\[
  C\;:=\;\{x\in X_\sigma:x(f)\le-2\alpha\}.
\]
Since $f$ is a fixed finite linear combination of flags and
$x\mapsto x(f)$ is a continuous map $X_\sigma\to\R$, the set $C$ is
closed.  We have just shown $P_n(C)\ge c_{k,\rho}$ for all large $n$.

\paragraph{Step 5: Pass to the limit.}
By \Cref{thm:razborov}(ii), the measures $P_n$ converge weakly to
$\Ext_\sigma(\phi_0)$ as a measure on $X_\sigma$.  By the closed-set
portmanteau theorem,
\[
  \Prob_{\phi^\sigma\sim\Ext_\sigma(\phi_0)}\!\bigl[\phi^\sigma\in C\bigr]
  \;\ge\;\limsup_{n\to\infty}P_n(C)
  \;\ge\; c_{k,\rho}\;>\;0.
\]
Hence $\Prob[\phi^\sigma(f)\le-2\alpha]>0$, contradicting
\ref{item:ensemble}.  This contradiction
proves~\ref{item:quotient}, and the theorem.
\qed

\section{Discussion}\label{sec:discussion}

\subsection{What clone-closure buys, what it costs}

\Cref{thm:main} should be read as: when $\K$ is closed under blow-ups,
labelled quotient semantics is no stronger than ensemble semantics
anchored at a $\K$-supported unlabelled root.  Practically, this means
any flag-algebra calculation on the constrained theory can be done
inside the unconstrained algebra $\A^\sigma[\T]$, using the random
extension formalism, with no loss of information.  The price is the
hypothesis on $\K$.

For $K_r$-free graphs the hypothesis is automatic and the result is
clean.  For more exotic hereditary classes one must check
clone-closure case by case; the criterion in \Cref{ex:Kr-free-clone-closed}
(forbidden patterns with no false twins) is sufficient but not
necessary.

\subsection{Comparison with the removal-lemma route}

For the unlabelled type ($\sigma=0$), the equivalence
\ref{item:quotient}$\,\Leftrightarrow\,$\ref{item:ensemble} holds for
\emph{every} hereditary $\K$, by a much shorter argument: the random
extension $\Ext_0(\phi_0)$ degenerates to the Dirac mass at $\phi_0$,
so~\ref{item:ensemble} becomes the statement that
$\phi_0(f)\ge 0$ for every $\K$-supported root, which is exactly
~\ref{item:quotient} pulled back through
\Cref{lem:quotient-support}.  In other words, the difficulty in the
present paper is specific to labels.  For $\sigma=0$ the random
extension is its own root, so the union of all ensembles already
exhausts the $\K$-supported homomorphisms on $\A^0[\T]$.  For
$\sigma\neq 0$, \Cref{lem:support-passes} keeps random extensions
$\K$-supported, but it leaves open whether the union of ensembles
$\bigcup_{\phi_0}\operatorname{supp}\Ext_\sigma(\phi_0)$ is rich
enough, as a sub-collection of $\K$-supported homomorphisms on
$\A^\sigma[\T]$, to detect every asymptotic inequality.  Our theorem
says: under clone-closure, it is---this collection's closed convex hull
already contains every $\K$-supported homomorphism (see
\Cref{cor:image}).

For $\sigma=0$ there is also a more constructive route via the
\emph{(induced) graph removal lemma}.  Let
$\psi\in\Homp(\A^0[\T_\K],\R)$ and write
$\widetilde\psi:=\psi\circ\q_0\in\Homp(\A^0[\T],\R)$ for its
$\K$-supported pull-back.  By Razborov's representation theorem
(\Cref{thm:razborov}(i)) applied in $\T$, some $\T$-sequence $\{G_n\}$
converges to $\widetilde\psi$.  Since $\widetilde\psi(H)=0$ for every
forbidden $H\in\mathcal H$ (where $\mathcal H$ defines $\K$), the
induced removal lemma of
Alon--Fischer--Krivelevich--Szegedy~\cite{AFKS2000} and its hereditary
generalisation by Alon--Shapira~\cite{AlonShapira2008} repair $G_n$ to
$G_n'\in\K$ within $o(n^2)$ edges, giving a $\K$-sequence with the same
limit.  This produces a concrete sequence of $\K$-models with the
desired limit without invoking Razborov's representation inside
$\T_\K$.

The blow-up proof of this paper is dual to both: it works for any
labelled type $\sigma$, at the cost of imposing clone-closure on $\K$.
The three regimes are summarised in the comparison table of \S1.4.
Currently no single argument covers the combination
(arbitrary hereditary $\K$, arbitrary type $\sigma$).

\subsection{What is needed to drop the hypothesis on \texorpdfstring{$\K$}{K}}

\begin{question}\label{q:general}
Does the equivalence \ref{item:quotient}\,$\Leftrightarrow$\,\ref{item:ensemble}
of \Cref{thm:main} hold for every hereditary $\K$ and every non-degenerate
type~$\sigma$ in $\T_\K$?
\end{question}

The blow-up proof fails because an independent blow-up of a $\K$-model
need not be in $\K$ (\Cref{rmk:why-hypothesis}).  Any approach replacing
the blow-up by a less destructive ``host'' construction would have to
satisfy two competing constraints: it must enlarge a finite labelled
flag $(G,\theta)\in\F^\sigma[\T_\K]$ into a sequence of unlabelled
$\K$-models inside which a random extension recovers $(G,\theta)$ with
positive probability, and it must preserve the densities of the relevant
flags.  In the unlabelled case, the removal lemma provides exactly such
a perturbation; what one wants in general is, roughly, a \emph{labelled}
induced removal lemma that respects a planted labelled embedding.

\subsection{Extensions and corollaries}

The abstract setup of \Cref{thm:main} subsumes several useful special
cases:

\begin{itemize}[leftmargin=2em]
\item \emph{Bipartite and multipartite graphs.}  The class of bipartite
graphs is the same as $\mathrm{Forb}^{\mathrm{ind}}(C_{2k+1}:k\ge 1)$,
and is clone-closed because an independent blow-up of a bipartite graph
inherits a bipartition (cluster-wise).  The same is true of the class
of $r$-partite graphs for every fixed $r$.
\item \emph{Linear hypergraphs.}  The class of $r$-uniform hypergraphs
in which every two hyperedges share at most one vertex is clone-closed:
two hyperedges of $B_N$ sharing two vertices project to two hyperedges
of $G$ sharing two vertices.
\item \emph{Sparse oriented graph classes.}  For every $g\ge 3$, the
class of oriented graphs of girth at least $g$ (i.e.\ with no
$\overrightarrow C_j$ for any $3\le j<g$) is clone-closed.  Suppose for
contradiction that a base oriented graph $G$ of girth $\ge g$ admits a
blow-up $B_N$ containing a $\overrightarrow C_\ell$ with $\ell<g$, and
take $\ell$ smallest.  The cyclic projection of this $\overrightarrow
C_\ell$ to $V(G)$ is a closed walk of length $\ell$ in $G$ with
consecutive vertices distinct (arcs in $B_N$ are inter-cluster).  Such a
closed walk contains a directed cycle of length between $3$ and $\ell$
(at most $\ell$, and at least $3$ because $G$ is oriented), contradicting
$\mathrm{girth}(G)\ge g>\ell$.  This includes the girth conditions
featured in the Caccetta--H\"aggkvist
conjecture~\cite[\S5]{Razborov2007}.

\smallskip
By contrast, ``$\overrightarrow C_k$-free for a single $k$'' is in
general \emph{not} clone-closed: if $G$ contains an
$\overrightarrow C_d$ with $d\ge 3$ a proper divisor of $k$, blowing up
$G$ with clusters of size $\ge 2$ traverses this short cycle $k/d$
times and produces a $\overrightarrow C_k$ absent from $G$.  Concretely,
blowing up a single $\overrightarrow C_3$ with all cluster sizes equal
to $2$ contains a $\overrightarrow C_6$, so the class of
$\overrightarrow C_6$-free oriented graphs is not clone-closed.  This is
why we state the result for full girth bounds, not for single
forbidden lengths.
\item \emph{Clique-free uniform hypergraphs.}  For every $r$-uniform
hypergraph theory and every fixed $s\ge r$, the class of hypergraphs
containing no induced complete $r$-uniform sub-hypergraph on $s$
vertices is clone-closed (a complete sub-hypergraph in the blow-up uses
at most one vertex per cluster and projects to one in the base).
\end{itemize}

\subsection{A clean restatement}

We close with an equivalent, geometric way to phrase \Cref{thm:main}:

\begin{corollary}[Quotient homomorphisms lie in the ensemble convex hull]\label{cor:image}
Under the hypotheses of \Cref{thm:main}, set
\[
  \mathcal{E}_{\K,\sigma}
  \;:=\;
  \bigcup_{\substack{\phi_0\ \K\text{-supported}\\\phi_0(\brak\sigma)>0}}
  \;\operatorname{supp}\Ext_\sigma(\phi_0)
  \;\subseteq\; \Homp(\A^\sigma[\T],\R).
\]
Then every $\psi\in\Homp(\A^\sigma[\T_\K],\R)$, viewed as a $\K$-supported
element of $\Homp(\A^\sigma[\T],\R)$ via \Cref{lem:quotient-support},
lies in the closed convex hull of $\mathcal{E}_{\K,\sigma}$ taken in the
product topology on $\R^{\F^\sigma}$.
\end{corollary}

\begin{proof}
Suppose not.  Then $\psi$ is separated from the closed convex hull by a
continuous linear functional on $\R^{\F^\sigma}$.  Such a functional
depends on only finitely many coordinates, so it is the evaluation
$x\mapsto x(g)$ for some $g\in\A^\sigma[\T]$, and the separation gives
$c\in\R$ with $\psi(g)<c$ while $\phi^\sigma(g)\ge c$ for every
$\phi^\sigma\in\mathcal{E}_{\K,\sigma}$.  Replacing $g$ by
$f:=g-c\cdot 1_\sigma$ converts this into
\[
  \psi(f)<0,
  \qquad
  \phi^\sigma(f)\ge 0\text{ for every }\phi^\sigma\in\mathcal{E}_{\K,\sigma}.
\]
The second condition is exactly $\Prob[\phi^\sigma(f)\ge 0]=1$ for every
ensemble, that is, condition \ref{item:ensemble} of \Cref{thm:main}.
The theorem therefore gives $\psi'(\q_\sigma f)\ge 0$ for every
$\psi'\in\Homp(\A^\sigma[\T_\K],\R)$, and applying this to the $\psi'$
whose pull-back is $\psi$ contradicts $\psi(f)<0$.
\end{proof}

This may be read as a partial structural answer, for blow-up-closed
classes, to the family of questions discussed in Razborov's
\S6~\cite[\S6]{Razborov2007} on the relationship between labelled
quotient homomorphisms and ensembles of random homomorphisms: under
clone-closure, every labelled quotient homomorphism is at least a
weak-limit of convex combinations of ensemble members, even if it need
not literally appear in the support of any single ensemble.

\begin{thebibliography}{9}

\bibitem{Razborov2007}
A.~A.~Razborov.
\newblock Flag algebras.
\newblock \emph{Journal of Symbolic Logic}, 72(4):1239--1282, 2007.

\bibitem{AFKS2000}
N.~Alon, E.~Fischer, M.~Krivelevich, and M.~Szegedy.
\newblock Efficient testing of large graphs.
\newblock \emph{Combinatorica}, 20(4):451--476, 2000.

\bibitem{AlonShapira2008}
N.~Alon and A.~Shapira.
\newblock A characterization of the (natural) graph properties testable
with one-sided error.
\newblock \emph{SIAM Journal on Computing}, 37(6):1703--1727, 2008.

\end{thebibliography}

\end{document}
